% ---- Encoding & language ----
\usepackage[utf8]{inputenc}
\usepackage[T1]{fontenc}
\usepackage[english]{babel}
\usepackage{lmodern}
\usepackage{microtype}

% ---- Math ----
\usepackage{amsmath,amssymb,amsthm,mathtools}
\usepackage{physics}      % \bra, \ket, \dd, \pdv, \comm, etc.
\usepackage{bm}
\usepackage{siunitx}

% ---- Graphics & layout ----
\usepackage{graphicx}
\usepackage[margin=2.7cm,headheight=14pt]{geometry}
\usepackage{xcolor}
\usepackage{tikz}
\usetikzlibrary{decorations.pathreplacing,arrows.meta,positioning,calc}
% Uncomment when needed:
% \usepackage{quantikz}                       % quantum circuit diagrams
% \usetikzlibrary{circuits.ee.IEC}            % electrical-engineering symbols
\usepackage{caption}
\usepackage{subcaption}
\usepackage{float}

% ---- Colour palette ----
% A small, deliberately restrained palette: a deep blue for the main
% structural elements, a teal accent for paragraph headings and
% definitions, a warm amber for examples, and a muted purple for
% remark-like environments. The numeric saturation is low everywhere
% so the document still reads as a quiet book, not a slide deck.
\definecolor{primary}{RGB}{30, 70, 130}
\definecolor{accent}{RGB}{0, 110, 125}
\definecolor{example}{RGB}{180, 105, 35}
\definecolor{remark}{RGB}{95, 55, 130}

% ---- Headers / footers ----
\usepackage{fancyhdr}
\pagestyle{fancy}
\fancyhf{}
\renewcommand{\headrulewidth}{0pt}
\fancyhead[LE]{\small\sffamily\color{primary!70}\thepage \hspace{1em} \leftmark}
\fancyhead[RO]{\small\sffamily\color{primary!70}\rightmark \hspace{1em} \thepage}
\renewcommand{\chaptermark}[1]{\markboth{\textsc{Ch.~\thechapter}\quad #1}{}}
\renewcommand{\sectionmark}[1]{\markright{\thesection\quad #1}}
\fancypagestyle{plain}{%
  \fancyhf{}\renewcommand{\headrulewidth}{0pt}%
  \fancyfoot[C]{\small\sffamily\color{primary!70}\thepage}%
}

% ---- Section / chapter / paragraph formatting ----
\usepackage{titlesec}

% Big numeral on the right, title to the right of it, thin rule below.
\titleformat{\chapter}[display]
  {\normalfont\sffamily\bfseries}
  {\filleft\color{primary!35}\fontsize{72}{72}\selectfont\thechapter}
  {-12pt}
  {\filleft\color{primary}\Huge}
  [\vspace{6pt}{\color{primary}\titlerule[1.2pt]}]
\titlespacing*{\chapter}{0pt}{-30pt}{30pt}

% Unnumbered chapters (Abstract, Introduction, Bibliography, ...) get
% the same look without the numeral.
\titleformat{name=\chapter,numberless}[display]
  {\normalfont\sffamily\bfseries}{}{0pt}
  {\filleft\color{primary}\Huge}
  [\vspace{6pt}{\color{primary}\titlerule[1.2pt]}]

\titleformat{\section}
  {\normalfont\sffamily\Large\bfseries\color{primary}}
  {\thesection}{0.7em}{}
\titlespacing*{\section}{0pt}{1.6\baselineskip}{0.5\baselineskip}

\titleformat{\subsection}
  {\normalfont\sffamily\large\bfseries\color{primary!85}}
  {\thesubsection}{0.7em}{}
\titlespacing*{\subsection}{0pt}{1.2\baselineskip}{0.4\baselineskip}

% Run-in paragraph headings, accent colour, terminated with a period.
\titleformat{\paragraph}[runin]
  {\normalfont\sffamily\bfseries\color{accent}}
  {}{0pt}{}[.]
\titlespacing*{\paragraph}{0pt}{0.6\baselineskip}{0.5em}

% Coloured part pages.
\titleformat{\part}[display]
  {\normalfont\sffamily\bfseries\color{primary}\centering}
  {\Large\MakeUppercase{\partname}~\thepart}{1em}{\Huge}

% ---- References & links ----
\usepackage{hyperref}
\hypersetup{
  colorlinks=true,
  linkcolor=primary,
  citecolor=accent,
  urlcolor=primary,
  pdfusetitle,
}
\usepackage{cleveref}

% ---- Theorem-like environments (numbering) ----
\theoremstyle{plain}
\newtheorem{theorem}{Theorem}[chapter]
\newtheorem{proposition}[theorem]{Proposition}
\newtheorem{lemma}[theorem]{Lemma}
\newtheorem{corollary}[theorem]{Corollary}

\theoremstyle{definition}
\newtheorem{definition}[theorem]{Definition}
\newtheorem{statement}[theorem]{Statement}
\newtheorem{postulate}[theorem]{Postulate}
\newtheorem{example}[theorem]{Example}
\newtheorem{application}[theorem]{Application}

\theoremstyle{remark}
\newtheorem*{phenomenon}{Phenomenon}
\newtheorem*{intuition}{Intuition}
\newtheorem*{remark}{Remark}
\newtheorem*{note}{Note}

% ---- Boxed presentation of the theorem-like environments ----
% mdframed wraps an existing amsthm environment with a coloured
% background and a thick left rule, without changing how the
% environment is numbered or how its label is referenced. We
% deliberately use the default framemethod (not TikZ) so the boxes
% break across pages cleanly when a long theorem outgrows a page.
\usepackage{mdframed}

\mdfdefinestyle{thmbox}{%
  rightline=false, topline=false, bottomline=false, leftline=true,
  linecolor=primary, linewidth=3pt,
  backgroundcolor=primary!4,
  innerleftmargin=10pt, innerrightmargin=10pt,
  innertopmargin=8pt, innerbottommargin=8pt,
  skipabove=10pt, skipbelow=10pt,
}
\mdfdefinestyle{defbox}{%
  rightline=false, topline=false, bottomline=false, leftline=true,
  linecolor=accent, linewidth=3pt,
  backgroundcolor=accent!4,
  innerleftmargin=10pt, innerrightmargin=10pt,
  innertopmargin=8pt, innerbottommargin=8pt,
  skipabove=10pt, skipbelow=10pt,
}
\mdfdefinestyle{exbox}{%
  rightline=false, topline=false, bottomline=false, leftline=true,
  linecolor=example, linewidth=3pt,
  backgroundcolor=example!4,
  innerleftmargin=10pt, innerrightmargin=10pt,
  innertopmargin=8pt, innerbottommargin=8pt,
  skipabove=10pt, skipbelow=10pt,
}
\mdfdefinestyle{remarkbox}{%
  rightline=false, topline=false, bottomline=false, leftline=true,
  linecolor=remark, linewidth=2pt,
  backgroundcolor=remark!3,
  innerleftmargin=8pt, innerrightmargin=8pt,
  innertopmargin=6pt, innerbottommargin=6pt,
  skipabove=8pt, skipbelow=8pt,
}

\surroundwithmdframed[style=thmbox]{theorem}
\surroundwithmdframed[style=thmbox]{proposition}
\surroundwithmdframed[style=thmbox]{lemma}
\surroundwithmdframed[style=thmbox]{corollary}
\surroundwithmdframed[style=thmbox]{statement}
\surroundwithmdframed[style=thmbox]{postulate}
\surroundwithmdframed[style=defbox]{definition}
\surroundwithmdframed[style=exbox]{example}
\surroundwithmdframed[style=exbox]{application}
\surroundwithmdframed[style=remarkbox]{phenomenon}
\surroundwithmdframed[style=remarkbox]{intuition}
\surroundwithmdframed[style=remarkbox]{remark}
\surroundwithmdframed[style=remarkbox]{note}

% ---- Custom shortcuts ----
% The physics package (and occasionally babel) define some of these
% names already, so we \let\name\undefined before each definition.
% That makes the redefinition unconditional whether or not the name
% was already in use.
\let\hilbert\undefined
\newcommand{\hilbert}{\mathcal{H}}
\let\id\undefined
\newcommand{\id}{\mathbf{1}}
\let\C\undefined
\newcommand{\C}{\mathbb{C}}
\let\R\undefined
\newcommand{\R}{\mathbb{R}}
\let\Z\undefined
\newcommand{\Z}{\mathbb{Z}}
\let\N\undefined
\newcommand{\N}{\mathbb{N}}
\let\hc\undefined
\newcommand{\hc}{\text{h.c.}}
\let\Tr\undefined
\DeclareMathOperator{\Tr}{Tr}
\let\op\undefined
\newcommand{\op}[1]{\hat{#1}}
